\documentclass[10pt,a4paper]{article}
\usepackage{amsmath,amssymb,amsthm,graphicx}
\usepackage{amsmath}
\usepackage{enumerate}
\usepackage{amsfonts}
\usepackage{algpseudocode}
\usepackage{algorithm}
\usepackage{amsthm}
\newtheorem{theorem}{Theorem}
\newtheorem{proposition}[theorem]{Proposition}
\newtheorem{lemma}[theorem]{Lemma}
\newtheorem{corollary}[theorem]{Corollary}
\newtheorem{definition}[theorem]{Definition}
\newtheorem{example}[theorem]{Example}
\newcommand{\out}{\textup{out}} 
\newcommand{\spa}{\textup{span}} 
\newcommand{\Stab}{\textup{Stab}} 
\newcommand{\peb}{\textup{peb}} 

\newcommand{\AcceptEdge}[1]{\textbf{Accept-Edge}(#1)}
\newcommand{\MovePebble}[1]{\textbf{Move-Pebble}(#1)}
\title{Bipartite (k1,k2,l)-sparsity is equivalent to arboricity of the associated butterfly multigraph after adding l-1 edges}
\author{Signe Lundqvist, Tovohery Randrianarisoa, Klara Stokes, Joannes Vermant}
\date{arxiv:2306.05050v3; CC BY 4.0}
\begin{document}
\maketitle
\noindent\textit{Editorial note.} In the source the theorem's own label is written after the statement text (source line 198); the retained span covers the whole theorem environment, label included.\par\medskip
\noindent\textit{Editorial note.} The two citations inside the retained lines (Servatius on k-plane matroids; Nash-Williams on edge-disjoint spanning trees) are resolved by a flattened two-entry bibliography fragment, so no bibliography database and no BibTeX run is needed.\par\medskip
\noindent\textit{Editorial note.} The theorem is stated in the source with k1, k2, l in the positive integers and k1 + k2 > l, although the surrounding discussion also treats l in the nonnegative integers; the statement is retained verbatim and no hypothesis was added.\par\medskip
\noindent\textit{Editorial note.} The retained environments print with this package's own numbering rather than the numbering of the article's full document.\par\medskip
\noindent\textit{Editorial scope.} One original exact result: the equivalence between (k1,k2,l)-sparsity of a bipartite graph and the forest decomposition of its butterfly multigraph after adding any l-1 edges.\par\medskip
\noindent\textit{Allowed standard prerequisites (editorial).}\begin{itemize}
\item Nash-Williams arboricity theorem. A finite multigraph decomposes into k edge-disjoint forests if and only if every nonempty edge subset F' satisfies |F'| <= k(|V(F')| - 1); cited by the author at source line 214.
\item Butterfly bookkeeping. The (k1,k2)-fold butterfly graph of a bipartite graph carries k1k2 copies of every edge, so an edge count for the bipartite graph is recovered by summing over the k1k2 pairs of vertex classes.
\end{itemize}\medskip
\begin{definition}\label{def: bipartite def}
    Let $G=(V_1, V_2, E)$ be a bipartite graph, with a partition of the vertex set into sets $V_1$ and $V_2$, and let $k_1, k_2 \in \mathbb{Z}_{>0}, l\in \mathbb{Z}_{\geq 0}$. Then we say that $G$ is $(k_1, k_2, l)$-sparse if, for every nonempty subset of $E' \subseteq E$, it holds that
\begin{equation*}
     \vert E'\vert \leq k_1 \vert V(E') \cap V_1 \vert +  k_2 \vert V(E')\cap V_2 \vert - l.
\end{equation*}
    We say that $G$ is $(k_1, k_2, l)$-tight if it is $(k_1, k_2, l)$-sparse, and if additionally
\begin{equation*}
     \vert E\vert = k_1 \vert V(E) \cap V_1 \vert +  k_2 \vert V(E)\cap V_2 \vert - l.
\end{equation*}
\end{definition}

To end this section, we will relate
$(k_1,k_2,l)$-sparsity to tree packings. The special case of $(k,1,k)$-tight and $(1,k,k)$-tight graphs was shown by Servatius \cite{k-planematroids}. 

Given a bipartite graph $G=(V_1,V_2,E)$, define its associated $(k_1,k_2)$-fold butterfly graph $B$ as follows: the vertex set consists of $k_1$ vertices for each element of $V_1$ and $k_2$ vertices for each element of $V_2$. For each edge $vw \in E$ with $v\in V_1$ and $w\in V_2$, there are edges between all $k_1$ vertices representing $v$ and all $k_2$ vertices representing $w$. 

\begin{theorem}
    Let $k_1, k_2, l\in \mathbb{Z}_{>0}$ with $k_1+k_2 >l$. A bipartite graph $G=(V_1,V_2,E)$ is $(k_1,k_2,l)$-sparse if and only if adding any $l-1$ edges between $V_1$ and $V_2$ results in a bipartite graph with an associated $(k_1,k_2)$-fold butterfly graph $B=(V, F)$ that decomposes into $k_1k_2$ edge-disjoint forests.
    \label{arboricity}
\end{theorem}

\begin{proof}
Suppose first that any resulting $B$ can be written as the union of $k_1k_2$ edge-disjoint forests. Then any such $B$ is $(k_1k_2, k_1k_2)$-sparse by the Nash-Williams theorem. For any nonempty subsets $V_1' \subseteq V_1$ and $V_2'\subseteq V_2$, adding the $(l-1)$ edges within $V_1' \cup V_2'$, one has, by the sparsity of the resulting graph $B$, that $$k_1k_2| E(V_1' \cup V_2')|+ (l -1)k_1k_2 \leq k_1k_2(k_1| V_1'|) + k_1k_2(k_1| V_1'|) - k_1k_2,$$ from which the $(k_1, k_2, l)$-sparsity of $G$ follows immediately.

 Conversely, suppose that we have added any $l-1$ edges to $G$, and let $B$ be the resulting butterfly graph. Let $V'\subseteq V(B)$. Note that all edges belonging to $F(V')$ are between $V_{1, i}'$ and $V_{2, j}'$ for some copies of subsets of $V_1$ and $V_2$ respectively. Since $G$ is sparse, we can bound $|F(V_{1, i}'\cup  V_{2, j}')|$ by sparsity:
    \begin{equation*}
        |F(V_{1, i}'\cup  V_{2, j}')| \leq k_1 |V_{1, i}'| + k_2 |V_{2, j}'| - 1
    \end{equation*}
    Then we have
    \begin{align*}
        |F(V')| &= \sum_{i=1}^{k_1} \sum_{j=1}^{k_2} |F(V_{1, i}'\cup  V_{2, j}')|\\
        &\leq  \sum_{i=1}^{k_1} \sum_{j=1}^{k_2}\left( k_1 |V_{1, i}'| + k_2 |V_{2, j}'| - 1\right)\\
        &=  k_2 k_1 \sum_{i=1}^{k_1} |V_{1, i}'| + k_1 k_2 \sum_{j=1}^{k_2} |V_{2, j}'| - k_1 k_2\\
        &=  k_2 k_1 (|V'| - 1).
    \end{align*}
    The result then follows by the Nash-Williams theorem \cite{Nash-Williams}. 
\end{proof}

\begin{thebibliography}{9}
\bibitem{Nash-Williams} C. St. J. A. Nash-Williams. Edge-disjoint spanning trees of finite graphs. \emph{J. London Math. Soc.} 36 (1961), 445--450.
\bibitem{k-planematroids} Brigitte Servatius. \emph{$k$-plane matroids and Whiteley's flattening conjectures}. In: Combinatorics, graph theory and computing, Springer Proc. Math. Stat. 388, Springer, Cham, 2022, 109--115.
\end{thebibliography}
\end{document}
